\documentclass[10pt]{article}
\usepackage[mono]{sim-oficial}
\input{datos-proponente}
\usepackage{ragged2e}
\newcommand{\clau}[2]{\par\needspace{4\baselineskip}\textbf{#1:} #2\par}
\newcommand{\encab}[2]{\begin{center}\bfseries #1\\[2pt] #2\end{center}}
\setlength{\parskip}{5pt}
\renewcommand{\arraystretch}{1.15}
\begin{document}
\begin{center}
\textbf{REPERTORIO N\textdegree{} SIM-2012-1.118}\\[6pt]
\textbf{CONSTITUCIÓN DE SOCIEDAD POR ACCIONES}\\
\textbf{LAFROX SpA}\\[6pt]
{\small Extracto para el Art.~39\textdegree, numeral 1, de las Bases Administrativas.\\Notaría, repertorios, inscripciones y personas son ficticios.}
\end{center}
\hrule\vspace{6pt}

En Valparaíso, República de Chile, a catorce de marzo de dos mil doce, ante mí, \textbf{\MakeUppercase{\SIMnotarioTitular}}, Abogado, Notario Público titular de la \SIMnotaria, con oficio en \SIMnotariaDir, comparecen: don \textbf{\SIMsocioAnombre}, chileno, soltero, \SIMprofesion, cédula nacional de identidad número \SIMsocioArun, domiciliado en \SIMsocioAdom; doña \textbf{\SIMsocioBnombre}, chilena, casada, ingeniera comercial, cédula nacional de identidad número \SIMsocioBrun, domiciliada en \SIMsocioBdom; don \textbf{\SIMsocioCnombre}, chileno, casado, ingeniero civil industrial, cédula nacional de identidad número \SIMsocioCrun, domiciliado en \SIMsocioCdom; y doña \textbf{\SIMsocioDnombre}, chilena, soltera, ingeniera en información y control de gestión, cédula nacional de identidad número \SIMsocioDrun, domiciliada en \SIMsocioDdom; todos mayores de edad, quienes acreditan su identidad con las cédulas citadas y exponen: que han acordado constituir, como en efecto constituyen, una sociedad por acciones regida por las normas de los artículos 424 y siguientes del Código de Comercio, introducidos por la Ley N\textdegree{} 20.190, y supletoriamente por la Ley N\textdegree{} 18.046 sobre Sociedades Anónimas, y por las estipulaciones siguientes:

\encab{TÍTULO PRIMERO}{ESTATUTOS}

\clau{PRIMERO}{La sociedad se denominará \textbf{``LafroX SpA''}, pudiendo usar el nombre de fantasía \textbf{``LafroX''}.}
\clau{SEGUNDO}{El domicilio de la sociedad será la comuna de Valparaíso, sin perjuicio de las agencias o sucursales que se establezcan en otros lugares del país o del extranjero.}
\clau{TERCERO}{La sociedad tendrá por objeto: a) el diseño, desarrollo, integración, comercialización y operación de sistemas y plataformas de software; b) la consultoría informática, la gestión de infraestructura tecnológica y la prestación de servicios de tecnologías de la información, incluida la mesa de servicios y el monitoreo de operaciones; c) la comercialización y soporte de equipamiento tecnológico; y d) la realización de toda otra actividad conexa o complementaria que acuerden los accionistas.}
\clau{CUARTO}{La duración de la sociedad será indefinida.}
\clau{QUINTO}{El capital de la sociedad es la suma de \SIMcapitalinicial{} (cien millones de pesos), dividido en 100.000 (cien mil) acciones nominativas, de una misma serie y sin valor nominal, que los accionistas suscriben y pagan íntegramente en dinero efectivo en este acto, en la forma siguiente:}

\begin{center}\small
\begin{tabular}{@{}l r r r@{}}
\hline
Accionista & Acciones & Participación & Aporte (\$)\\\hline
\SIMsocioAnombre & 40.000 & 40 \% & 40.000.000\\
\SIMsocioBnombre & 25.000 & 25 \% & 25.000.000\\
\SIMsocioCnombre & 20.000 & 20 \% & 20.000.000\\
\SIMsocioDnombre & 15.000 & 15 \% & 15.000.000\\\hline
\textbf{Total} & \textbf{100.000} & \textbf{100 \%} & \textbf{100.000.000}\\\hline
\end{tabular}
\end{center}

\clau{SEXTO}{La sociedad será administrada por un Directorio de tres miembros y un Gerente General. Se designa Presidenta del Directorio a doña \SIMpresidenta, cédula de identidad número \SIMpresidentarun. Se designa Gerente General a don \SIMgerente, cédula de identidad número \SIMgerenterun, quien será el representante legal de la sociedad y la representará judicial y extrajudicialmente con las facultades de un mandatario con libre administración y las que le otorgue el Directorio, pudiendo delegar parte de ellas en apoderados.}
\clau{SÉPTIMO}{Los accionistas se reunirán en junta ordinaria dentro del primer cuatrimestre de cada año y en junta extraordinaria cuando lo exijan las necesidades sociales. Los acuerdos se adoptarán por mayoría absoluta de las acciones emitidas con derecho a voto, salvo los casos en que la ley o estos estatutos exijan un quórum superior.}
\clau{OCTAVO}{La sociedad confeccionará un balance general al treinta y uno de diciembre de cada año. Salvo acuerdo unánime en contrario, se distribuirá como dividendo al menos el treinta por ciento de las utilidades líquidas.}
\clau{NOVENO}{La sociedad se disolverá por las causales legales. Las diferencias entre los accionistas, o entre éstos y la sociedad, se someterán a arbitraje de derecho ante árbitro designado de común acuerdo o, a falta de éste, por la justicia ordinaria de Valparaíso.}

Un extracto de esta escritura se inscribió a fojas \SIMcbrfojas{} número \SIMcbrnumero{} del Registro de Comercio de \SIMcbrciudad{} del año \SIMcbranio{}, con fecha 21 de marzo de 2012, y se publicó en el Diario Oficial el 28 de marzo de 2012, edición N\textdegree{} SIM-40.501.

En comprobante y previa lectura, firman los comparecientes. Se da copia. Doy fe.

\vspace{6pt}
\noindent\begin{tabularx}{\linewidth}{@{}X X@{}}
\rule{5.6cm}{0.4pt}\newline\small\SIMsocioAnombre\newline RUN \SIMsocioArun &
\rule{5.6cm}{0.4pt}\newline\small\SIMsocioBnombre\newline RUN \SIMsocioBrun\\[16pt]
\rule{5.6cm}{0.4pt}\newline\small\SIMsocioCnombre\newline RUN \SIMsocioCrun &
\rule{5.6cm}{0.4pt}\newline\small\SIMsocioDnombre\newline RUN \SIMsocioDrun\\
\end{tabularx}

\needspace{10\baselineskip}
\encab{PRIMERA MODIFICACIÓN}{REPERTORIO N\textdegree{} SIM-2018-3.402 --- AUMENTO DE CAPITAL Y DIVISIÓN DE ACCIONES}

Por escritura pública de veintidós de agosto de dos mil dieciocho, otorgada en la \SIMnotaria, la junta extraordinaria de accionistas acordó aumentar el capital social de \SIMcapitalinicial{} a \SIMcapital{} (mil millones de pesos), mediante la capitalización de utilidades retenidas por \$ 400.000.000 y el aporte en dinero de los accionistas, a prorrata, por \$ 500.000.000. Asimismo, se acordó una \textbf{división de acciones en razón de uno a diez}, de modo que cada una de las 100.000 acciones existentes quedó dividida en diez acciones, sin alterar la participación de ningún accionista. En consecuencia, el artículo QUINTO de los estatutos queda así: ``El capital de la sociedad es de \SIMcapital, dividido en \SIMaccionestotal{} de acciones nominativas, de una misma serie y sin valor nominal''. Se inscribió al margen de la inscripción matriz, a fojas 3.502 número 1.644 del año 2018, y su extracto se publicó en el Diario Oficial el 12 de septiembre de 2018, edición N\textdegree{} SIM-42.109.

\encab{SEGUNDA MODIFICACIÓN}{REPERTORIO N\textdegree{} SIM-2019-2.775 --- DOMICILIO DE OPERACIONES}

Por escritura pública de diez de septiembre de dos mil diecinueve, otorgada en la \SIMnotaria, la junta extraordinaria acordó reemplazar el artículo SEGUNDO de los estatutos por el siguiente: ``El domicilio de la sociedad será la comuna de Valparaíso, y su centro de operaciones funcionará en Av. Brasil 2241, Valparaíso, sin perjuicio de su sede central y de las agencias o sucursales que establezca en otros lugares''. Se inscribió a fojas 5.117 número 2.286 del año 2019, y su extracto se publicó en el Diario Oficial el 26 de septiembre de 2019, edición N\textdegree{} SIM-42.457.

\textbf{Estructura de capital vigente:}
\begin{center}\small
\begin{tabular}{@{}l r r r@{}}
\hline
Accionista & Acciones & Participación & Capital (\$)\\\hline
\SIMsocioAnombre & \SIMsocioAacc & \SIMsocioApct{} \% & 400.000.000\\
\SIMsocioBnombre & \SIMsocioBacc & \SIMsocioBpct{} \% & 250.000.000\\
\SIMsocioCnombre & \SIMsocioCacc & \SIMsocioCpct{} \% & 200.000.000\\
\SIMsocioDnombre & \SIMsocioDacc & \SIMsocioDpct{} \% & 150.000.000\\\hline
\textbf{Total} & \textbf{\SIMaccionestotal} & \textbf{100 \%} & \textbf{1.000.000.000}\\\hline
\end{tabular}
\end{center}

\needspace{9\baselineskip}
\textbf{AUTORIZACIÓN.} Certifico que el presente extracto concuerda con las escrituras públicas originales de constitución y modificaciones de LafroX SpA, que he tenido a la vista. Valparaíso, veintitrés de noviembre de dos mil veintiséis. Copia autorizada simulada con fines académicos. Doy fe.

\vspace{4pt}
\noindent\begin{minipage}[b]{0.4\linewidth}\SIMtimbrenotarial[3cm]\end{minipage}\hfill
\begin{minipage}[b]{0.5\linewidth}\centering\rule{6cm}{0.4pt}\\\SIMnotarioTitular\\Notario Público\\\SIMnotaria\end{minipage}
\end{document}
